\section{Dark-Sector Hypotheses and Current Constraints}
\subsection{A load-dependent vacuum density}
The proposed interpretation identifies vacuum energy with substrate maintenance or thermal noise. A changing load does not by itself specify a density law or an equation of state. If one additionally posits $\rho(a)=\rho_*\mathcal L(a)$ and the separately conserved homogeneous continuity equation, then
\begin{equation}
a\rho'(a)=-3(1+w)\rho(a)
\end{equation}
gives
\begin{equation}
w(a)=-1-\frac{a\mathcal L'(a)}{3\mathcal L(a)}.
\end{equation}
Lean proves the derivative and algebraic equivalence for nonzero density normalization and load. For $a>0$, this expression equals $-1$ exactly when $\mathcal L'(a)=0$. The sign and magnitude of any deviation require an evolution law for $\mathcal L$, and an interacting sector would require an additional exchange term in the continuity equation. No redshift-dependent equation of state has yet been selected by the substrate.

\CatchFileBetweenTags{\OmegaRatioVal}{\PMcalculations}{OmegaRatioVal}
\CatchFileBetweenTags{\OmegaRatioExperimentalValue}{\PMcalculations}{OmegaRatioExperimentalValue}
\CatchFileBetweenTags{\OmegaRatioSigma}{\PMcalculations}{OmegaRatioSigma}

\subsection{Mode counts are not yet energy-density ratios}
One abundance ansatz uses the Lie-algebra dimension $248$ and three assumed generations of $16$ active labels:
\begin{equation}
\frac{248}{3\nu}=\frac{31}{6}.
\end{equation}
This Lie-algebra dimension is distinct from the rank-eight lattice dimension. Its interpretation as a cosmological capacity count requires an explicit model. If the 48 active labels are a subset of 248 total labels and the dark component is their disjoint complement, the inactive-to-active count is instead
\begin{equation}
\frac{248-48}{48}=\frac{25}{6}.
\end{equation}
Lean checks this distinction. Moreover, density ratios require the energy weights and occupancies of the two components; equal counts alone do not fix those weights.

The calculator evaluates the corrected count expression
\begin{equation}
R_{\rm ansatz}=\frac{248}{3\nu(1-\chi/\Delta+\alpha)}
\approx\OmegaRatioVal.
\end{equation}
It differs from the draft reference \OmegaRatioExperimentalValue\ by approximately $\OmegaRatioSigma\sigma$. The correction terms and identification with $\Omega_{DM}/\Omega_b$ remain assumptions; this numerical proximity does not establish a dark-matter component, its clustering or the absence of detectable dark-sector particles.

\CatchFileBetweenTags{\MnuRestVal}{\PMcalculations}{MnuRestVal}
\CatchFileBetweenTags{\MnuWeakVal}{\PMcalculations}{MnuWeakVal}

\subsection{The earlier contextual-neutrino proposal does not fit the splittings}
The proposed contextual model uses a three-state inverse-resonance tower with largest rest mass near $9.209$ meV, total rest sum approximately $\MnuRestVal$ meV, and a common interaction multiplier six, giving total interaction sum approximately $\MnuWeakVal$ meV. These are candidate formulas, not an established distinction between gravitational and propagation masses.

The upstream Lean results now rule out the common-scale tower as a fit to the two fixed oscillation windows used in the audit. For masses proportional to $(1,1/43,1/43^2)$, the positive squared-mass gaps are proportional to
\begin{equation}
(1,1849,1850).
\end{equation}
A common multiplier cannot change their ratios, and no relabeling supplies the moderate hierarchy of the two comparison windows. Reading the atmospheric gap instead as $(6m_{\rm rest})^2-m_{\rm rest}^2$ exceeds the fixed Paper V benchmark by more than $17\sigma$. Paper VI imports both exclusions rather than repeating the earlier claim of a resolution.

A small rest-mass sum is therefore insufficient. A replacement must derive a propagation Hamiltonian and its eigenvalue splittings, establish how the gravitational and propagation readouts are related, and evaluate the cosmological inference under that dynamics. A nonuniform response law is not excluded by the common-scale obstruction, but has not been selected by the present substrate model.
